\section{Logarithmic compactification}
\label{sec:compactification}

The wall structure of Theorem~\ref{thm:annular-consistency} is proper on the
open annulus, but its supports have genuine exits at the deleted boundary and
focus ends.  Restoring the compact boundary does not turn these exits into
ordinary harmless boundary joints.  We replace them by logarithmic caps.
The construction is first made for a finite GKT window
$\Wcal=(L,P,\mathcal E)$ in the sense of
Definition~\ref{def:intro-finite-window}, a deformation order $N$, and a
finite root subtree $T$.  Its coefficient ring is
$R_{\Wcal,N}$ from \eqref{eq:intro-window-ring}, with $L$ in place of $I$.
The symbol $\Xfrak_{T,\Wcal,N}$ denotes the finite logarithmic cap obtained
from the root subtree $T$ over this ring.

\subsection{Why the ends must be modified}

Near a smooth boundary point choose affine coordinates in which the tangent
wedge is the upper half-plane and the boundary tangent is $(1,0)$.  If a
wall terminates transversely with primitive inward direction $m=(u,v)$,
$v>0$, and outgoing multiplicity $d>0$, the extremal bend changes the
boundary tangent to
\begin{equation}\label{eq:boundary-bend}
 (1+dvu,dv^2).
\end{equation}
Its normal component is positive.  The extremal broken line remains inside
the wedge, so the boundary-convexity argument for tangent canonical walls
does not apply \cite{grossetal2016theta}.  The endpoint retains nontrivial
one-sided transport.

Nor can one remove the boundary by an affine reflection.  In a boundary
tangent--normal basis, write its nontrivial holonomy and a putative
pointwise-fixing reflection as
\[
 U=\begin{pmatrix}1&w\\0&1\end{pmatrix},
 \qquad R=\begin{pmatrix}1&c\\0&-1\end{pmatrix}.
\]
The relation $RU=UR$ forces $2w=0$, whereas the two annular boundary widths
are $4$ and $16$.

At a finite stage, close the wall skeleton in $\overline\Acal$ and retain the
two one-sided germs at every transverse endpoint.  A path may terminate at
such a germ but is not continued across it.  Let
\begin{equation}\label{eq:exit-word}
 E_{\Wcal,N}=\Theta_m\cdots\Theta_1
\end{equation}
be the complete ordered product along one attaching circle, after coincident
abelian supports have been combined.  At $A,C$ this word is empty; at
$B_0,B_1$ all diagonal exits lie in one abelian group.  This stopped object
is a useful intermediate form, but it is not yet a logarithmic
compactification.

The focus germs already have standard Kummer nearby charts.  If $s$ is the
slab divisor and $xy=ts$ is the node chart, the degree-$d$ root is the fs
monoid map
\begin{equation}\label{eq:focus-monoid}
 \langle x,y,t,s\mid x+y=t+s\rangle
 \longrightarrow
 \langle x,y,t,r\mid x+y=t+dr\rangle,
 \qquad s\longmapsto dr.
\end{equation}
On the atlas this is $xy=tr^d$; inverting $r$ recovers the root-slab chart of
Section~\ref{sec:affine-kummer}.

\begin{lemma}[The focus root chart]\label{lem:focus-root-chart}
For every $d\geq1$, the monoid
\[
 P_d=\langle x,y,t,r\mid x+y=t+dr\rangle
\]
is fine and saturated.  Its monoid algebra is
\[
 \Bbbk[P_d]\simeq
 \Bbbk[x,y,t,r]/(xy-tr^d).
\]
The map in \eqref{eq:focus-monoid} is Kummer with group cokernel
$\ZZ/d$, and its atlas is finite free, given by the monic equation
$r^d=s$.  After any coefficient base change, both the chart and its
boundary $r=0$ are flat over the coefficient ring.
\end{lemma}

\begin{proof}
The defining binomial is prime because its exponent relation is primitive.
The hypersurface $xy=tr^d$ is Cohen--Macaulay.  For $d=1$ its singular
locus is the origin; for $d>1$ it is contained in $x=y=r=0$, of codimension
two.  Serre's criterion therefore makes $\Bbbk[P_d]$ normal, equivalently
$P_d$ is saturated.  The group cokernel and the monic presentation follow
from $s\mapsto dr$.  Finally,
\[
 \Bbbk[P_d]/(r)\simeq\Bbbk[x,y,t]/(xy),
\]
so tensoring either ring with a coefficient algebra over $\Bbbk$ preserves
flatness.
\end{proof}

\subsection{Affine cylinder ends and peripheral holonomy}

Real-oriented blowup replaces the two boundary components and four focus
points by six attaching circles.  Their shear widths are
\begin{equation}\label{eq:six-widths}
 (w_0,w_1,w_A,w_{B,0},w_{B,1},w_C)=(4,16,1,1,1,1).
\end{equation}

\begin{proposition}[Outward affine flares]\label{prop:outward-flares}
Let an attaching circle have affine length $\ell>0$ and holonomy $T^{-w}$.
The quotient of $r\geq0$ by
\begin{equation}\label{eq:flare-action}
 (x,r)\longmapsto(x+\ell+wr,r)
\end{equation}
is a flaring integral-affine cylinder.  The action is free and proper, the
circle at height $r$ has length $\ell+wr$, and $r$ is a proper integral
convex exhaustion.  Every finite-stage wall exit extends to a proper ray and
cannot return to the annular core.
\end{proposition}

\begin{proof}
The $n$th power of \eqref{eq:flare-action} translates $x$ by
$n(\ell+wr)$.  This quantity is nonzero for $n\ne0$, which proves freeness
and properness.  The normal component of a tangent vector is preserved by
the shear.  An exiting wall with positive normal component therefore crosses
each level of $r$ at most once.
\end{proof}

The affine cylinder has the correct affine holonomy, but its peripheral loop
also crosses the product \eqref{eq:exit-word}.  Let $\rho(H)$ denote the
prescribed affine, root, sector, and volume-line transport around the
attaching circle.
Choose a radial seam disjoint from the exit supports and glue its two sides by
\begin{equation}\label{eq:corrected-seam}
 G_{\Wcal,N}=\rho(H)E_{\Wcal,N}^{-1}.
\end{equation}
A positive peripheral loop crosses $E_{\Wcal,N}$ and then the seam, so its
total transport remains
\begin{equation}\label{eq:seam-identity}
 G_{\Wcal,N}E_{\Wcal,N}=\rho(H).
\end{equation}
The stopped operation has not been killed; it is the correction between the
bare peripheral gluing and the cap seam.

The polarization extends over every flare with positive integral kink seven.
Indeed, if one branch on the core is $ax+br+c$, use
$ax+(b+7)r+c$ on the flare.  The value seven is a convenient uniform choice:
the directions and pole orders of the four basic wall characters at their
negative ends are
\[
\begin{array}{c|rrrr}
 q&(2,1)&(4,3)&(3,4)&(1,2)\\ \hline
 Lq&(-4,3)&(-8,7)&(-6,7)&(-2,3)\\
 \text{pole order}&3&7&7&3.
\end{array}
\]
Additivity gives a bound of seven per deformation degree.  Equivalently, the
auxiliary Rees relation
\begin{equation}\label{eq:Rees-relation}
 \lambda=u^7s
\end{equation}
regularizes every finite ordered product of exiting-wall automorphisms in one
dominating expansion.  Equation~\eqref{eq:Rees-relation} does not define the
boundary: its divisor $u=0$ is
vertical over $\lambda=0$.  Horizontal flatness will instead come from
transporting the monoid chart itself.

\begin{proposition}[Separation of the exits]\label{prop:exit-separation}
At each finite stage, the positive and negative support families on either
outer flare are pairwise disjoint, apart from coincident supports combined in
one abelian ray group.  A support-free seam can therefore be chosen, and the
chamber nerve of the cut flare is an interval.  The focus flares $A,C$ carry
no support, while $B_0,B_1$ carry one abelian radial support.
\end{proposition}

\begin{proof}
Develop an outer attaching circle in invariant tangent--normal coordinates.
A support of tangential speed $\sigma$ has equation
$x(r)=\phi_i-\sigma(r+h_i)$, while the deck period at height $r$ is
$w_i(r+h_i)$.  Thus two supports meet modulo deck translation exactly when
their slope difference lies in $w_i\ZZ$.  The exact slope intervals, derived
in Appendix~\ref{app:calculations}, have width smaller than $w_i$; every
positive-minus-negative difference lies strictly between $0$ and $w_i$.
Hence no two distinct supports meet.  The focus assertions follow from the
cofactor placement: only the diagonal ray reaches a $B$ lift.
\end{proof}

\subsection{Transported monoids}

Let $R=R_{\Wcal,N}$ be the finite Artin coefficient ring just defined.  Choose a base
chamber $C_0$, an fs toric cap monoid $P$, and a finite Kummer algebra $S_0$
over $R[P]$.  Locally, $S_0$ is obtained by tangent localizations and monic
root equations
\begin{equation}\label{eq:monic-roots}
 S_i=S_{i-1}[r_i]/(r_i^{d_i}-f_i).
\end{equation}
Let $\nu:P\to\NN$ be the outward valuation and
\begin{equation}\label{eq:boundary-ideal}
 J_0=(z^p:\nu(p)>0)S_0.
\end{equation}
If $\Phi_C$ is the ordered Jacobi product from $C_0$ to a chamber $C$, set
inside the punctured Kummer algebra
\begin{equation}\label{eq:transported-charts}
 S_C=\Phi_C(S_0),
 \qquad J_C=\Phi_C(J_0),
 \qquad P\longrightarrow S_C,
 \quad p\longmapsto\Phi_C(z^p).
\end{equation}
The transport acts on the monoid chart itself.  A character
with a pole in $S_0$ is regular in the appropriate $S_C$ by construction.

Around a wall endpoint choose a small toric open star $U$.  If $\Phi_-$ and
$\Phi_+$ are the products in the adjacent chambers, identify the complete
transported charts, including their boundary strata, by
\begin{equation}\label{eq:full-bridge}
 \Phi_+\Phi_-^{-1}:
 \Phi_-(S_0(U))\xrightarrow{\sim}\Phi_+(S_0(U)).
\end{equation}
The identification includes the full boundary overlap.  Gluing only after
inverting the outward uniformizer would leave two copies of the boundary
point.

\begin{theorem}[Finite-stage logarithmic cap]
\label{thm:finite-log-cap}
The charts \eqref{eq:transported-charts}, glued by the full boundary identifications
\eqref{eq:full-bridge}, form a separated fine-saturated formal logarithmic
cap over $R$.  Its asymptotic divisor is defined chamberwise by $J_C$ and is
flat over $R$.  The complement of this divisor is the corrected mapping
torus \eqref{eq:corrected-seam}; reduction modulo the deformation ideal is
the ordinary toric/Kummer cap with peripheral transport $\rho(H)$.  The
construction commutes with restriction of the root tree and labelled coefficient system
and with reduction of $N$.
\end{theorem}

\begin{proof}
Every equation in \eqref{eq:monic-roots} is monic.  Before quotienting by
the finite diagonalizable root group, $S_0$ is therefore finite free over
the preceding toric ring; after restriction to the toric boundary the same
monic equations remain.  Since $R$ is a characteristic-zero $\Bbbk$-algebra,
the Reynolds operator splits invariants under this finite group.  Thus the
invariant affine charts, as well as the root-stack atlases themselves, have
$R$-flat boundary, and in particular $S_0/J_0$ is $R$-flat.

Transport by $\Phi_C$ is an $R$-algebra isomorphism.  The prelog chart in
\eqref{eq:transported-charts} is strictly isomorphic to the original toric
chart, so its characteristic monoid is fs, and
\begin{equation}\label{eq:horizontal-flatness}
 S_C/J_C\simeq S_0/J_0
\end{equation}
is $R$-flat.

By Proposition~\ref{prop:exit-separation}, the wall closures meet the
asymptotic divisor at distinct points and the chamber nerve is an interval.
The bridge \eqref{eq:full-bridge} is an isomorphism of the entire chart
rings.  On endpoint bridges the overlap ring is either chart ring; on
chamber-band overlaps, conjugation by $\Phi_C$ reduces to an ordinary
intersection of toric opens.  In each case the affine-intersection
multiplication map is surjective, proving separatedness.  Since the endpoint
stars are disjoint, there are no further intersections or cap joints.

If $C_m$ is the last chamber, $\Phi_{C_m}=E_{\Wcal,N}$.  At the seam,
\[
 G_{\Wcal,N}(S_{C_m})
 =\rho(H)E_{\Wcal,N}^{-1}E_{\Wcal,N}(S_0)
 =\rho(H)(S_0),
\]
which closes the cap with total transport \eqref{eq:seam-identity}.  Modulo
the deformation ideal every wall product is the identity.  All operations
are functorial under the stated restrictions.
\end{proof}

\begin{figure}[t]
\centering
\vspace{0.5\baselineskip}
\begin{adjustbox}{max width=\textwidth,center}
\begin{tikzpicture}[
  x=0.72cm,y=0.72cm,
  every node/.style={font=\small},
  box/.style={draw,rounded corners=2pt,fill=blue!4,minimum width=1.25cm,
    minimum height=.75cm,align=center},
  arr/.style={-{Latex[length=2mm]},line width=.7pt},
  wall/.style={line width=1pt,blue!70!black},
  boundary/.style={line width=1.1pt,orange!80!black}]
  \path[use as bounding box] (-0.5,-2.25) rectangle (15.65,3.0);

  % Panel a: developed affine flare.
  \begin{scope}[shift={(0,0)}]
    \node[font=\small\bfseries] at (2.25,2.65) {(a) Outward flare};
    \draw[boundary] (0,0) -- (4.5,0)
      node[midway,below=3pt] {attaching circle};
    \draw[gray!60] (0.45,2.05) -- (4.05,2.05);
    \draw[wall] (0.65,0) -- (1.12,2.05);
    \draw[wall] (1.55,0) -- (2.20,2.05);
    \draw[wall] (2.65,0) -- (3.52,2.05);
    \draw[dashed,gray!70] (4.2,0) -- (4.2,2.05)
      node[midway,left=2pt] {seam};
    \draw[arr] (-0.15,.25) -- (-0.15,1.75)
      node[midway,left=1pt] {$r$};
    \node[align=center] at (2.25,-1.05)
      {period $\ell+wr$\\chamber nerve: interval};
  \end{scope}

  % Panel b: transported monoid charts and full bridges.
  \begin{scope}[shift={(5.25,0)}]
    \node[font=\small\bfseries] at (2.25,2.65) {(b) Full bridges};
    \node[box,minimum width=1.05cm] (s0) at (.45,1.0) {$S_0$\\$J_0$};
    \node[box] (s1) at (2.25,1.0) {$S_1$\\$J_1$};
    \node[box,minimum width=1.05cm] (s2) at (4.05,1.0) {$S_2$\\$J_2$};
    \draw[arr] (s0.east) --
      node[above=5pt,font=\scriptsize,fill=white,inner sep=.5pt] {$\Phi_1$}
      (s1.west);
    \draw[arr] (s1.east) --
      node[above=5pt,font=\scriptsize,fill=white,inner sep=.5pt] {$\Phi_{21}$}
      (s2.west);
    \node at (2.25,1.95) {$P\longrightarrow S_C$};
    \draw[boundary] (0.05,-.10) -- (4.45,-.10);
    \draw[fill=orange!10,draw=orange!75!black] (1.45,-.42) rectangle (3.05,.22);
    \node[font=\footnotesize] at (2.25,-.10) {full bridge};
    \node[orange!80!black] at (2.25,-.72) {horizontal divisor $J_C$};
    \node[font=\footnotesize] at (2.25,-1.25)
      {$S_C=\Phi_C(S_0)$};
    \node[font=\footnotesize] at (2.25,-1.62)
      {$S_C/J_C\simeq S_0/J_0$};
  \end{scope}

  % Panel c: Rees domination versus horizontal flatness.
  \begin{scope}[shift={(10.65,0)}]
    \node[font=\small\bfseries] at (2.3,2.65) {(c) Rees versus horizontal};
    \node[box,minimum width=2.6cm] (rees) at (2.3,1.55)
      {Rees chart\\$\lambda=u^7s$};
    \node[red!70!black,fill=white,inner sep=1pt] at (2.3,.72)
      {vertical: $u=0\Longrightarrow\lambda=0$};
    \node[box,minimum width=2.9cm] (hor) at (2.3,-.35)
      {horizontal cap\\$S_C/J_C$ flat over $R$};
    \draw[arr] (rees.east) -| ++(.55,0) |- node[right=2pt,pos=.72]
      {dominates} (hor.east);
    \node[align=center] at (2.3,-1.35)
      {coefficients independent of $u$};
  \end{scope}
\end{tikzpicture}

\end{adjustbox}
\caption{A logarithmic cap.  Wall rays on the outward flare give a linear
chamber order.  Each chamber carries the transported monoid chart
$P\to S_C$, and adjacent charts are identified across a full bridge
containing the horizontal boundary.  The Rees relation controls pole order
but does not define that boundary.}
\label{fig:log-cap}
\end{figure}

There are three distinct restrictions of a finite cap.  Removing its
asymptotic divisor gives the punctured corrected mapping torus.  Reducing the
coefficient ring modulo its deformation ideal gives the toric/Kummer
scattering-special fiber.  Restricting to the asymptotic divisor gives
$\Spf(S_C/J_C)$, which remains flat over the entire coefficient ring by
\eqref{eq:horizontal-flatness}.  No coefficient variable is identified with
a spatial boundary uniformizer.

\subsection{The compactification theorem}

Truncate each flare at an integral level.  The extension of the polarization
has positive kink seven, and the product subdivision gives locally convex
boundary.  The original focus points have become peripheral end loops; their
primitive monodromies remain as holonomy rather than codimension-two
singularities of the truncation.  The finite tropical spaces are therefore
positive, locally rigid, and pre-polarized in the sense used in
\cite{carletal2024}.

\begin{theorem}[Separated fs logarithmic compactification]
\label{thm:regular-compactification}
The finite caps of Theorem~\ref{thm:finite-log-cap}, including the focus
charts \eqref{eq:focus-monoid}, form an inverse system indexed by finite root
subtrees, finite labelled coefficient systems closed under the dependencies
of Appendix~\ref{app:gkt-completion}, and deformation orders.
Each stage is separated and fs logarithmic, and its six horizontal boundary
divisors are flat over its finite Artin coefficient ring.  Its affine base
has flare widths $4,16,1,1,1,1$ and a positive integral polarization.  Its
punctured restriction carries the original annular Jacobi wall structure,
with all root, sector, exponent, and volume-line holonomies unchanged.

The inverse limit and chamber-groupoid quotient give a logarithmic
compactification in the adic proalgebraic pro-Kummer category.  No
regularity or log-regularity assertion is made, and the limit is not claimed
to be an algebraic stack of finite type.
\end{theorem}

\begin{proof}
Proposition~\ref{prop:outward-flares} supplies the affine ends and their
proper convex exhaustion.  Proposition~\ref{prop:exit-separation} supplies a
support-free seam and linearly ordered chamber cover, so every end satisfies
Theorem~\ref{thm:finite-log-cap}.  Tensoring the Kummer monoid charts
\eqref{eq:focus-monoid} with the cap monoids extends the root data.
Lemma~\ref{lem:focus-root-chart} gives fs charts and coefficient-flat
boundary at these four ends; root stacks commute with the base change.
Proposition~\ref{prop:Jacobi-transport}
then extends the volume line and its bracket.

The full bridges introduce neither doubled boundary points nor new joints.
All new charts restrict to the original ones before the exits, so
contractible consistency and every framed annular product remain unchanged.
Monic root equations and \eqref{eq:horizontal-flatness} give flatness at
each stage.  Root construction, ordered wall products, and bridge maps commute
with every finite restriction.  Passing to the compatible inverse system and
then to the groupoid quotient proves the theorem.
\end{proof}

The special fiber in this theorem is the toric/Kummer cap of the annular
model.  It is not identified with the central fiber of the original
disk-shaped quartic degeneration.  This distinction is topological as well
as categorical, and the topological half of it is an obstruction.

\begin{proposition}[No ordinary CPS recovery]
\label{prop:no-CPS-degeneration}
Neither the annular cover $\overline\Acal$ nor its noncompact
Legendre-dual cylindrical form is the dual intersection complex of a
distinguished toric degeneration of del Pezzo surfaces with simple
singularities, proper generic fiber, and smooth anticanonical divisor.
\end{proposition}

\begin{proof}
For such a degeneration the dual intersection complex is homeomorphic to
$\RR^2$ \cite[Prop.~7.13]{carletal2024}, so its first cohomology vanishes.
An annulus and an open cylinder both have $H^1(-,\Bbbk)\simeq\Bbbk$.
\end{proof}

The conclusion is limited to that class of degenerations.  It does not
exclude logarithmic or orbifold degenerations, non-proper models, a boundary
of different topology, or a comparison which agrees only on part of the
sector grading.  It does exclude reading the annulus as an unnoticed
substitute for the ordinary Carl--Pumperla--Siebert affine base.

\subsection{A regular dominating presentation}

The finite stages of Theorem~\ref{thm:regular-compactification} are not
themselves log regular in general.  This failure comes from the transported
root divisors, not from the asymptotic boundary.

\begin{proposition}[A multiple-root obstruction]
\label{prop:root-stage-nonregular}
There is a finite root-stack presentation whose atlas is singular and not log regular.
More precisely, three divisors in the orbit \eqref{eq:root-tower} meet with
distinct tangent directions at one point of the two-dimensional torus.
\end{proposition}

\begin{proof}
Write $f_A=1+x_A$ and
\[
 F_n=1+f_A^{4n}x_B.
\]
At $p=(x_A,x_B)=(-2,-1)$ one has $f_A(p)=-1$ and $F_n(p)=0$ for every
$n$.  For $n=0,1,2$ the gradients are
\[
 \nabla F_0(p)=(0,1),\qquad
 \nabla F_1(p)=(4,1),\qquad
 \nabla F_2(p)=(8,1),
\]
so the three reduced divisors are smooth and pairwise transverse at $p$.
The simultaneous root atlas has equations
\[
 r_i^2=F_i,\qquad 0\leq i\leq2.
\]
It is finite over a surface.  At the point above $p$ with all $r_i=0$, its
three-by-five Jacobian has rank two, so its tangent space has dimension
three.  The atlas is therefore singular.  If the three root divisors are put
in the log structure, the log ideal cuts out the closed point but the sharp
characteristic group has rank three.  Kato's dimension equality would read
$2=0+3$, so the chart is not log regular.  Such a finite stage occurs by
taking a rooted subtree containing the paths which give $F_0,F_1,F_2$.
\end{proof}

Regularity is recovered by a modification of the finite presentations.  The
distinction is useful: it keeps the minimal root object of
Theorem~\ref{thm:regular-compactification}, while providing a toroidal model
when regular geometry is required.

\begin{theorem}[Regular dominating presentation]
\label{thm:regular-presentation}
For every finite stage $\Xfrak_{T,\Wcal,N}$ of
Theorem~\ref{thm:regular-compactification}, over its Artin coefficient ring
$R=R_{\Wcal,N}$, there is a cofiltered system of proper morphisms
\[
 \Xfrak^{\mathrm{reg}}_{\lambda}
 \longrightarrow \Xfrak_{T,\Wcal,N}
\]
with the following properties.

\begin{enumerate}
\item Each morphism is an isomorphism away from the spatial boundary and the
transported slab-zero divisors.

\item Each $\Xfrak^{\mathrm{reg}}_{\lambda}$ is a separated logarithmic
Deligne--Mumford stack, log smooth over $(\Spec R,0)$, and has a smooth atlas
over $R$.  Every geometric fiber has a regular atlas and is log regular.

\item Its total logarithmic boundary is simple normal crossings on an
etale atlas.  In particular, the six asymptotic divisors remain flat over
$R$.

\item All finite root phases, the volume line, the Jacobi bracket,
the wall maps, and the corrected peripheral transports pull back.  Full
bridges and seams lift to strict logarithmic isomorphisms.

\item Common refinements are compatible with enlargement of the root tree,
restriction of $\Wcal$, and reduction of $N$.  Their inverse system carries
the chamber-groupoid action and gives a regularized adic pro-Kummer
compactification dominating that of
Theorem~\ref{thm:regular-compactification}.
\end{enumerate}

For a fixed resolved boundary, its component-root maps are Kummer log
etale.  A transition which introduces a previously absent root divisor is
not asserted to be Kummer log etale.
\end{theorem}

\begin{proof}
Fix the finite stage and form the reduced union $\Delta_T$ of its toric
boundary and its finitely many transported root divisors.  The root
radicands have deformation degree zero, so this pair is defined over
$\Bbbk$ and then base-changed to $R$.  First apply functorial logarithmic
resolution in characteristic zero to the root-free toroidal chart:
\[
 \rho:Z_T^{\mathrm{sm}}\longrightarrow Z_T.
\]
This is an isomorphism off $\Delta_T$, and $Z_T^{\mathrm{sm}}$ is a smooth
toroidal orbifold.  On it, functorially principalize the product of the
pulled-back radicand ideals, obtaining
\[
 \pi:Y_T\longrightarrow Z_T^{\mathrm{sm}},
 \qquad \widetilde\pi=\rho\circ\pi:Y_T\longrightarrow Z_T.
\]
Then $Y_T$ is smooth and the reduced total transform $D_T$ is simple normal
crossings \cite{abramovicheta2020principalization}.  Both steps are
functorial for toroidal isomorphisms, hence for the chamber maps, full
bridges, finite diagonal actions, and seams.

Etale locally, let $x_1,\ldots,x_s$ cut out the components of $D_T$.  For
every radicand $f_\alpha$ of degree $d_\alpha\in\{2,4\}$, principalization
gives
\[
 \widetilde\pi^*f_\alpha
 =u_\alpha\prod_{j=1}^s x_j^{a_{\alpha j}},
 \qquad u_\alpha\in\Ocal_{Y_T}^{\times}.
\]
Take the fourth root stack of every component, writing $x_j=y_j^4$, and
etale locally retain all solutions of
$\eta_\alpha^{d_\alpha}=u_\alpha$.  Then
\begin{equation}\label{eq:resolved-root-formula}
 g_\alpha=\eta_\alpha
 \prod_{j=1}^s y_j^{4a_{\alpha j}/d_\alpha}
 \qquad\text{satisfies}\qquad
 g_\alpha^{d_\alpha}=\widetilde\pi^*f_\alpha.
\end{equation}
All exponents are integral.  Thus the component-root stack maps to the
original root-stack presentation and retains its full root-choice orbit.  At a $B$ cusp,
$f_B=h^2$, formula \eqref{eq:resolved-root-formula} reads
$h=y^4$ and $g=\eta y^2$.  The ambient $\mu_4$ action exchanges the two
normalized quadratic branches, while the effective class on either branch
has order two, as in \eqref{eq:cusp-roots}.

Before the finite diagonal quotient, an etale local atlas is smooth over
$R$ with boundary $y_1\cdots y_s=0$.  Its chart is
$0\to\NN^s$, so Kato's chart criterion proves log smoothness
\cite{kato1989log}.  Finite diagonal groups are etale in characteristic zero;
their quotients are therefore log-smooth Deligne--Mumford stacks.  Geometric
fibers have regular atlases with simple-normal-crossings boundary and hence
are log regular.  The fourth-root map
$\NN^s\to\frac14\NN^s$ is Kummer log etale.  Smoothness over $R$ and the
normal-crossings description also give flatness of the asymptotic divisors.

The resolution--principalization composite is an isomorphism on the punctured
transport locus, so all wall, line, and theta data pull back unchanged.  Given
two finite regularizations, pass to the rooted union of their stages, take the
closure of their common graph, and apply the same
resolution--principalization procedure.  This gives a regular model
dominating both.  The resulting index category is cofiltered.  Path
concatenation reindexes it, while functoriality under isomorphisms lifts the
groupoid action.  This proves the pro-system assertion.  Finally, adding a
new branch divisor raises the characteristic rank at its generic point, so
such an enlargement map need not be Kummer; this explains the qualifier in
the statement.
\end{proof}
